\section{Chương~\ref{sec:chemoutput}. Đầu ra hóa học}

Hai dây dẫn tới tim và tốc độ khác nhau của chúng, phản hồi âm trong các chuỗi tầng hormone, mã bằng tần số các đợt, bộ tạo nhịp ngày đêm và việc chỉnh nó bằng ánh sáng đều đã được đo trực tiếp; giới hạn $K<8$ là một định lý của lý thuyết điều khiển, được suy ra trong chương; các xung sinh ra bởi chính phản hồi của chuỗi tầng là giả thuyết có bằng chứng thực nghiệm mạnh ủng hộ.

{\footnotesize
\setlength{\tabcolsep}{3pt}\renewcommand{\arraystretch}{1.15}
\begin{longtable}{>{\raggedright}p{0.5cm}>{\raggedright}p{4.0cm}>{\raggedright}p{5.6cm}>{\raggedright}p{2.3cm}>{\raggedright\arraybackslash}p{1.7cm}}
\toprule
TT & Khẳng định & Thí nghiệm và điều đã đo & Nguồn & Trạng thái \\
\midrule\endfirsthead
\toprule
TT & Khẳng định & Thí nghiệm và điều đã đo & Nguồn & Trạng thái \\
\midrule\endhead
1 & Bộ tạo nhịp tim tự đập khoảng $100$ lần mỗi phút; lúc nghỉ phanh đang được đạp, và tần số thường vào khoảng $60$--$70$ & Người, cắt hoàn toàn cả hai dây dẫn tới tim: tần số riêng cao hơn tần số lúc nghỉ và giảm theo tuổi, ở người trưởng thành khoảng $100$ nhịp mỗi phút. Vậy lúc nghỉ phanh chiếm ưu thế. Tần số nghỉ $60$--$70$ là giá trị điển hình, không trích dẫn riêng. & \cite{JoseCollison1970ihr} & đã đo \\
2 & Ga \nt{NORA} và phanh \nt{ACET} là các bộ lọc thông thấp, phanh nhanh hơn~\eqref{eq:gasbrake} & Chó, kích thích điều tần dây thần kinh phế vị hoặc dây giao cảm của tim và hàm truyền tới tần số tâm nhĩ: cả hai đường đều là bộ lọc thông thấp, đường giao cảm có tần số cắt thấp hơn nhiều và thêm một độ trễ thuần khoảng $1.7\,\text{s}$. Các đặc tính phụ thuộc vào trương lực trung bình, nên công thức tuyến tính~\eqref{eq:gasbrake} chỉ là gần đúng. & \cite{Berger1989} & đã đo \\
3 & Phanh nhanh vì \nt{ACET} mở kênh kali không cần chất truyền tin thứ hai, còn \nt{NORA} thì qua một chuỗi phản ứng & Mảnh màng tế bào tâm nhĩ: kênh kali do acetylcholine mở được mở bởi tiểu đơn vị $\beta\gamma$ của protein G gắn với thụ thể, không cần chất truyền tin thứ hai bên trong tế bào. Đường \nt{NORA} qua cAMP không được trích dẫn ở đây. & \cite{Logothetis1987gbg} & đã đo \\
4 & Máu đi hết một vòng cơ thể trong khoảng một phút; \nt{ADRE} từ dòng máu tới được mọi cơ quan có thụ thể & Ước tính chung về bậc độ lớn; trong chương cũng được phát biểu như vậy, không trích dẫn nguồn. & --- & ước tính \\
5 & Chuỗi tầng hormone được bao bởi phản hồi âm: hormone cuối ức chế các tầng đầu & Tổng quan thí nghiệm: hormone vỏ thượng thận ức chế tuyến yên tiết ACTH và vùng dưới đồi tiết hormone giải phóng trong ba dải thời gian: nhanh, trung gian và chậm. & \cite{KellerWood1984fb} & đã đo \\
6 & Ba tầng giống nhau ổn định khi $K<8$~\eqref{eq:cascadestable}, tại biên chu kỳ là $2\pi\tau/\sqrt3\approx3.6\tau$ & Suy ra trong chương: ở tần số $\sqrt3/\tau$, ba tầng cho lệch pha $180^\circ$ và suy giảm $8$ lần. & suy ra trong chương & lý thuyết \\
7 & Sai số của mức là $1/(1+K)$, nên bộ điều chỉnh như vậy không giữ chính xác hơn khoảng $10\,\%$ (mệnh đề~\ref{prop:cascade}) & Hệ quả của dòng~6 với phản hồi tỉ lệ. Trục thật có phản hồi ở nhiều tầng và trong nhiều dải thời gian (dòng~5), nên giới hạn này thuộc về mô hình~\eqref{eq:cascade}, chứ không phải đã đo. & suy ra trong chương & lý thuyết \\
8 & Khi $K=4$ và $7$ mức ổn định, khi $K=12$ có các xung đều với chu kỳ $3.7$--$3.9\,\tau$ (hình~\ref{fig:cascade}) & Tính theo \texttt{models/\allowbreak chem\_\allowbreak models.py}: khi $K=12$ các chu kỳ liên tiếp là $3.9$, $3.7$, $3.8$, $3.7\,\tau$; khi $K=4$ biên độ ở cuối phép tính là $0.003$. & \texttt{models/\allowbreak chem\_\allowbreak models.py} & mô hình \\
9 & Hormone stress được tiết theo xung khoảng một giờ một lần; các xung do chính phản hồi sinh ra, không cần bộ tạo nhịp riêng & Chuột cống, truyền liên tục hormone giải phóng: ACTH và corticosterone dao động với tần số sinh lý, gần một giờ một lần; không cần đầu vào nhịp nhàng từ vùng dưới đồi để có xung. Mô hình tuyến yên--thượng thận với phản hồi có trễ cho ra các dao động như vậy. & \cite{Walker2012glc, Walker2010} & giả thuyết \\
10 & Bên nhận đọc tần số các đợt chứ không đọc mức & Khỉ không tự tiết hormone giải phóng gonadotropin: truyền liên tục không phục hồi sự tiết hormone tuyến yên, còn các đợt mỗi giờ một lần thì phục hồi; chuyển sang truyền liên tục lại dập tắt đáp ứng. Thụ thể mệt đi như bộ lọc thông cao là cách diễn giải. & \cite{Belchetz1978} & đã đo \\
11 & Các tần số đợt khác nhau tác động khác nhau lên các loại tế bào khác nhau trong cùng một tuyến & Cùng những con khỉ đó: tăng tần số đợt lên $2$--$5$ lần mỗi giờ làm giảm cả hai hormone tuyến yên; giảm xuống một lần mỗi $3$~giờ làm giảm một hormone (LH) và tăng hormone kia (FSH), nên tỉ lệ giữa chúng do tần số quyết định. & \cite{Wildt1981gnrh} & đã đo \\
12 & Nhịp ngày đêm sinh ra từ phản hồi âm nội bào với độ trễ vài giờ & Tổng quan: protein của các gen đồng hồ ức chế có trễ sự phiên mã của chính các gen đó; vòng phiên mã--dịch mã cho chu kỳ khoảng một ngày. & \cite{TakahashiJS2017} & đã đo \\
13 & Bộ tạo nhịp chính là một nhóm vài chục nghìn nơ-ron, đồng bộ phần còn lại của cơ thể & Tổng quan: nhân trên giao thoa thị giác là bộ tạo nhịp ngày đêm chính; từng nơ-ron của nó khi nuôi cấy riêng là những bộ tạo nhịp độc lập, mạng lưới đồng bộ chúng; ở chuột nhắt có khoảng $10^4$ nơ-ron mỗi bên. & \cite{Welsh2010scn} & đã đo \\
14 & Chu kỳ riêng ở người là $24.18$~giờ & Người sống trong ánh sáng mờ theo lịch tách khỏi ngày đêm: chu kỳ của nhịp melatonin, nhiệt độ và cortisol trung bình $24.18$~giờ ở cả người trẻ và người già, với độ tản mát hẹp. & \cite{Czeisler1999} & đã đo \\
15 & Ánh sáng chỉnh bộ tạo nhịp như một vòng khóa pha~\eqref{eq:pll}; cần dịch khoảng $11$~phút mỗi ngày & Người, một lần chiếu sáng mạnh $6.7$~giờ vào các thời điểm khác nhau trong ngày: đường cong đáp ứng pha loại một với biên độ $5$~giờ, làm trễ pha trước điểm cực tiểu nhiệt độ cơ thể và làm sớm pha sau đó. Phương trình~\eqref{eq:pll} là mô hình chuẩn; $11$~phút $=0.18$~giờ là phép tính số học. & \cite{Khalsa2003prc} & đã đo \\
16 & Không có ánh sáng thì không có khóa pha, và nhịp trôi theo chu kỳ riêng & Người mù hoàn toàn, không cảm nhận ánh sáng, sống trong xã hội bình thường: khoảng một nửa có nhịp melatonin và cortisol chạy theo chu kỳ riêng, không trùng với ngày đêm. & \cite{Sack1992blind} & đã đo \\
\bottomrule
\end{longtable}}
